\section{Related work}\label{sec:related}
\subsection{Learned locomotion and actuator models}\label{sec:rw-loco}
Locomotion policies trained in simulation have been transferred to a real quadruped \cite{hwangbo2019agile}, and an accurate actuator model with simulated latency narrowed the gap between simulation and hardware \cite{tan2018simtoreal}. Shin et al.\ include the actuator torque-speed operating region in the training of a running quadruped \cite{shin2023actuator}, a limit of each actuator on its own. Bjelonic et al.\ add an energy model of permanent-magnet synchronous motors to sim-to-real training and state that it does not include battery dynamics \cite{bjelonic2026pace}. Tao et al.\ model voltage coupling and magnetic saturation of the actuators of a fast quadruped, with a fixed bus voltage and a voltage limit for each motor \cite{tao2026actuator}, and Zhang et al.\ restrict commanded joint accelerations to those that each joint's motor can realise under its voltage limit, with the DC-link voltage as a given bound and its ripple treated as model discrepancy \cite{zhang2026vra}. Miller et al.\ tune simulation parameters of a commercial quadruped against hardware data and attribute a remaining difference in swing-leg motion to power distribution limits that their model does not include \cite{miller2025spot}; those limits are our subject. Where these works state a joint's supply, it is ideal, fixed or a given bound; in our plant the bus voltage falls with the total load of both drives (Section~\ref{sec:model}).

\subsection{Energy, battery and thermal limits in legged robots}\label{sec:rw-energy}
Energy-aware methods reduce what a robot draws over a motion or a run. Valsecchi et al.\ \cite{valsecchi2024power} and Nai et al.\ \cite{nai2025power} learn the power drawn on hardware from real data and train policies that draw less; ECO (Energy-Constrained Optimization) turns energy consumption into a constraint enforced by the Lagrangian method during training \cite{huang2026eco}; and the RAIBO2 quadruped reduces losses in its mechanics, motor driving circuit and policy together, for range on one battery charge \cite{lee2026raibo2}. Battery models with a load-dependent voltage exist for robots: Evagora builds one with transient voltage sag for a high-power humanoid and leaves its coupling to robot control to future work \cite{evagora2025battery}, and Alcibar et al.\ predict the end-of-discharge voltage of inspection-drone batteries under varying load \cite{alcibar2025drone}. Shu et al.\ distribute power between a battery and a supercapacitor in a legged robot with learning-based model predictive control \cite{shu2025hess}, on the source side. Wan et al.\ \cite{wan2026thermal} and Qian et al.\ \cite{qian2026thermal} bring motor temperature into locomotion learning, through a residual policy and through thermal-constraint rewards. None that we found keeps the instantaneous bus voltage and battery current within limits during locomotion; we couple a battery model with voltage sag to a learned controller and enforce both limits on the load side at every substep.

\subsection{Battery and DC-bus safety in power electronics}\label{sec:rw-power}
Equivalent circuit models describe a battery's terminal voltage under load; Hu et al.\ compare twelve lumped models of this kind on dynamic test data \cite{hu2012ecm}, and our battery is a Thevenin model with one RC pair (Section~\ref{sec:model}). In DC microgrids, a decentralised controller based on control barrier functions (CBFs) keeps transient voltages and currents within bounds with an unknown or time-varying load \cite{kosaraju2022microgrid}, and a dynamic CBF with disturbance observers restrains the current of a storage unit under sudden load changes while regulating the DC bus voltage \cite{fu2026bess}. Safe reinforcement learning is also used in power and energy systems, from grid-connected converters to distribution networks \cite{bui2025review}: projection-based constrained policy optimisation reduces voltage, line and transformer limit violations in grid-to-vehicle management \cite{zhang2025g2v}, and constraint tightening with a projection layer gives constraint satisfaction of a learned microgrid controller \cite{liu2026ampc}. Schenke et al.\ guard the online training of a motor's torque controller with an algorithm that rejects unsafe switching actions, such as those that cause overcurrent \cite{schenke2023safe}. In these works the load is unknown or a disturbance, and the one learning-based electrical safeguard protects a single drive. Our filter acts on the load itself: the currents that several drives draw from one battery and bus.

\subsection{Safe reinforcement learning, safety layers and barrier functions}\label{sec:rw-safe}
Safe learning in robotics and safe reinforcement learning are reviewed in \cite{brunke2022safe,gu2024review}. Constrained methods act through training: constrained policy optimisation gives near-constraint satisfaction of expected costs at each policy iteration \cite{achiam2017cpo}. Stooke et al.\ show that the traditional Lagrange multiplier update behaves as integral control and add proportional and derivative terms \cite{stooke2020pid}; our SAC-Lagrangian baseline uses the traditional update. For legged robots, Kim et al.\ train locomotion with rewards and constraints to reduce reward engineering \cite{kim2024constraints}, and CaT turns constraint violations into stochastic terminations during policy learning \cite{chanesane2024cat}, which shapes training but does not act on the commands at execution, as our filter does. Safety layers act on the actions instead. Dalal et al.\ add a layer that corrects actions in closed form per state with a linearised model learned from data \cite{dalal2018safety}; RAYEN imposes hard convex constraints on the outputs of a neural network \cite{tordesillas2023rayen}; Cheng et al.\ combine reinforcement learning with CBF controllers and Gaussian-process dynamics to keep safety with high probability during learning \cite{cheng2019endtoend}; and CBF-RL filters rollouts with CBFs during training so that the policy can be deployed without an online filter \cite{yang2026cbfrl}. In legged locomotion, a learner policy can be switched to a safe recovery policy \cite{yang2022saferl}, or a recovery policy governed by a learned reach-avoid value network can avoid collisions \cite{he2024abs}. CBFs are an established tool for optimisation-based safety-critical control \cite{ames2019cbf}, and their discrete-time form \cite{agrawal2017dcbf} gives the form of our barrier conditions. CBF quadratic programs can become infeasible with many constraints and tight control limits \cite{xiao2022hocbf}, and unmodelled actuation disturbances degrade CBF guarantees, which input-to-state safety makes robust \cite{alan2023issf}. Our two conditions reduce to one power budget in closed form, so the filter knows at each substep whether currents meeting it exist, and errors in the battery parameters are met with margins and an EKF (Sections~\ref{sec:estimator}, \ref{sec:estresults} and~\ref{sec:stress}). Unlike the task-defined constraints above, ours come from the shared supply: the filter rebuilds them at every substep from the bus state and battery parameters, acts on the drive currents below the policy and stays in the loop in training and at deployment.

\subsection{Research gap and our response}\label{sec:gap}
Table~\ref{tab:related} compares these works on six features, using only what each paper states. Three gaps remain.
\begin{enumerate}
\item[G1.] Learned locomotion models the actuators in growing detail but treats each joint's supply as ideal, fixed or a given bus voltage \cite{bjelonic2026pace,tao2026actuator,zhang2026vra}; we found no training plant in which a shared battery and DC bus sag with the total load of all joints.
\item[G2.] Barrier-based safety for DC buses exists in power electronics, but it keeps bus voltages and currents, or a storage unit's current, within bounds while treating the load as unknown or as a disturbance \cite{kosaraju2022microgrid,fu2026bess}. Safety layers for reinforcement learning protect constraints on actions, collisions or expected costs \cite{dalal2018safety,tordesillas2023rayen,he2024abs,achiam2017cpo}, and the one electrical case we found guards a single drive rather than a shared supply \cite{schenke2023safe}.
\item[G3.] Energy-aware training and expected-cost constraints reduce energy or violations on average \cite{valsecchi2024power,nai2025power,huang2026eco,achiam2017cpo}, and a bound on an average does not rule out a breach in a single substep.
\end{enumerate}
We address G1 with the coupled battery, bus and drive model of Section~\ref{sec:model} (contribution~1); G2 with a load-side barrier filter whose budget is exact for one bus step (Proposition~\ref{prop:budget}) and which, with exact battery parameters, guarantees both limits for any policy whenever currents meeting the budget exist (Corollary~\ref{cor:closed}; contribution~2); and G3 with this per-substep guarantee, compared with reward penalties, a Lagrangian method and the same filter applied only at test time (Sections~\ref{sec:results} and~\ref{sec:ablation}; contribution~3).

\begin{table}[tbp]\centering\caption{Features of related works on actuator, power and constraint handling in learned control and DC microgrids, as stated in each paper: ``yes'' and ``no'' are stated in the paper, ``n.s.'' means not stated. Battery dynamics: a model of the battery's voltage (open-circuit voltage, internal resistance, sag). Shared bus: several drives on one DC bus whose voltage varies with their total load. Per-joint voltage limit: each motor's voltage limit (torque at speed limited by the motor's back electromotive force). Safety guarantee: a theorem that the constraint holds. Used in training: the model, constraint or filter acts during policy training. Estimation error: evaluated with errors in model parameters or estimates.}\label{tab:related}
{\footnotesize\setlength{\tabcolsep}{4pt}
\resizebox{\linewidth}{!}{\begin{tabular}{lcccccc}\toprule
Work & \begin{tabular}[b]{@{}c@{}}Battery\\dynamics\end{tabular} & \begin{tabular}[b]{@{}c@{}}Shared\\bus\end{tabular} & \begin{tabular}[b]{@{}c@{}}Per-joint\\voltage limit\end{tabular} & \begin{tabular}[b]{@{}c@{}}Safety\\guarantee\end{tabular} & \begin{tabular}[b]{@{}c@{}}Used in\\training\end{tabular} & \begin{tabular}[b]{@{}c@{}}Estimation\\error\end{tabular} \\\midrule
Shin et al.\ \cite{shin2023actuator} & n.s. & n.s. & yes & n.s. & yes & n.s. \\
Bjelonic et al.\ \cite{bjelonic2026pace} & no & n.s. & yes$^{a}$ & n.s. & yes & n.s. \\
Miller et al.\ \cite{miller2025spot} & n.s. & no$^{b}$ & n.s.$^{c}$ & n.s. & yes & n.s. \\
Nai et al.\ \cite{nai2025power} & n.s. & n.s. & n.s. & n.s. & yes$^{d}$ & n.s. \\
Huang et al.\ (ECO) \cite{huang2026eco} & n.s. & n.s. & n.s. & n.s. & yes & n.s. \\
Tao et al.\ \cite{tao2026actuator} & n.s. & n.s. & yes & n.s. & yes & n.s. \\
Zhang et al.\ \cite{zhang2026vra} & n.s. & n.s. & yes & no$^{e}$ & no$^{f}$ & yes$^{g}$ \\
Wan et al.\ \cite{wan2026thermal} & n.s. & n.s. & n.s. & n.s. & yes & n.s. \\
Qian et al.\ \cite{qian2026thermal} & n.s. & n.s. & n.s. & n.s. & yes & n.s. \\
Yang et al.\ (CBF-RL) \cite{yang2026cbfrl} & n.s. & n.s. & n.s. & yes$^{h}$ & yes & partly$^{i}$ \\
Tordesillas et al.\ (RAYEN) \cite{tordesillas2023rayen} & n.s. & n.s. & n.s. & yes & yes & n.s. \\
Chane-Sane et al.\ (CaT) \cite{chanesane2024cat} & n.s. & n.s. & n.s. & n.s. & yes & n.s. \\
Dalal et al.\ \cite{dalal2018safety} & n.s. & n.s. & n.s. & n.s. & yes & n.s. \\
Schenke et al.\ \cite{schenke2023safe} & n.s. & n.s.$^{j}$ & yes & n.s. & yes & yes$^{k}$ \\
Kosaraju et al.\ \cite{kosaraju2022microgrid} & n.s. & n.s.$^{l}$ & n.s. & n.s. & no & n.s. \\
\midrule
This work & yes & yes & yes & yes$^{m}$ & yes & yes$^{n}$ \\
\bottomrule\end{tabular}}}
\par\smallskip\parbox{\linewidth}{\footnotesize
$^{a}$~A torque envelope limited by the back electromotive force is enforced in simulation.
$^{b}$~Power distribution limits are stated not to be included in the model.
$^{c}$~Torque--speed limits are identified per joint; a voltage limit is not stated.
$^{d}$~A power model learned from hardware data is used as a reward.
$^{e}$~Proves that the commanded accelerations are realisable under the voltage limit, and states that this is not a safety guarantee.
$^{f}$~Applied at execution; the learned controller does not know the joint limits.
$^{g}$~Experiments across motors, temperatures and bus voltages.
$^{h}$~A bound on the barrier function of the discretised system, with a discretisation residual.
$^{i}$~Training with noise on the dynamics model.
$^{j}$~One drive; its DC-link voltage is measured and not assumed constant.
$^{k}$~Small violations are attributed to the prediction uncertainty of the online-identified drive model.
$^{l}$~A DC microgrid with an unknown or time-varying load; drives are not stated.
$^{m}$~For one substep with an exact battery model, under assumptions (A1)--(A7) (Corollary~\ref{cor:closed}).
$^{n}$~Errors in the battery's electromotive force and series resistance, a scaled plant resistance and an extended Kalman filter (Sections~\ref{sec:stress} and~\ref{sec:estresults}).}
\end{table}
